\documentclass[10pt,twocolumn]{article}

\usepackage[margin=0.75in]{geometry}
\usepackage[T1]{fontenc}
\usepackage[utf8]{inputenc}
\usepackage{times}
\usepackage{microtype}
\usepackage{amsmath,amssymb}
\usepackage{booktabs}
\usepackage{multirow}
\usepackage{xcolor}
\usepackage{graphicx}
\usepackage{url}
\usepackage[numbers,sort&compress]{natbib}
\usepackage[colorlinks=true,linkcolor=blue,citecolor=blue,urlcolor=blue]{hyperref}

\setlength{\columnsep}{0.22in}
\setlength{\parskip}{0pt}
\setlength{\parindent}{1em}

\title{\textbf{Disentangled Cross-Cultural Music Recommendation with\\Optimal Transport Alignment and Participatory Feedback}}
\author{Anonymous ISMIR 2026 Submission}
\date{}

\begin{document}
\maketitle

\begin{center}
\fcolorbox{red}{white}{\parbox{0.96\columnwidth}{\small
\textbf{Internal drafting note.}
This manuscript is a planning draft only.
All quantitative results marked with $\dagger$ are \textbf{synthetic placeholder values} introduced to stabilize the paper narrative before the final experiments are complete.
They must be replaced by real measurements before any circulation or submission.
}}
\end{center}

\begin{abstract}
Cross-cultural music recommendation remains difficult even with strong audio foundation models because stylistic and cultural variation is often entangled with affective or functional similarity.
As a result, retrieval systems may incorrectly treat unfamiliar music as irrelevant and collapse to popular or culturally dominant targets.
We present a proof-of-concept framework for cross-cultural recommendation that combines disentangled downstream representation learning, domain-adversarial alignment, and optimal transport (OT) ranking on top of foundation audio embeddings.
Our model factorizes each track representation into content, style, and relatively culture-agnostic affective-functional variables, then performs cross-cultural matching in the shared latent space instead of directly in the raw embedding space.
To correct high-uncertainty boundary cases, we further incorporate a participatory active learning (PAL) loop that converts a small amount of expert feedback into pairwise constraints for retraining.
In an internal draft evaluation on a four-domain cross-cultural dataset pipeline built from public audio sources, the proposed framework yields consistent placeholder gains over embedding-only baselines in ranking quality, serendipity, and cultural calibration.
We position this work as a runnable research prototype and an argument for combining foundation embeddings with structured downstream alignment in culturally sensitive recommendation settings.
\end{abstract}

\section{Introduction}
Recent music foundation models have improved audio representation learning for tagging, retrieval, and classification \citep{li2023mert,kanatas2025culturemert}.
However, stronger embeddings do not automatically solve cross-cultural recommendation.
In this setting, users often seek music that is unfamiliar in style yet meaningful in function, affect, or listening context.
Embedding-only retrieval pipelines can fail because stylistic unfamiliarity is frequently encoded as semantic distance, causing the system to recommend only the safest and most culturally dominant candidates.

We argue that cross-cultural recommendation should not be treated as standard nearest-neighbor retrieval in a monolithic embedding space.
Instead, it should be framed as a representation problem:
which parts of a music representation encode content, which encode stylistic-cultural variation, and which encode the relatively shared affective or functional structure needed to bridge across cultures?
This framing is motivated by recent cross-cultural MIR work showing both the importance of culturally adapted representations \citep{kanatas2025culturemert} and the need for better culturally grounded evaluation resources \citep{lee2025globalmood}.

Based on this view, we build a downstream model on top of foundation embeddings that separates three factors:
\emph{content} ($z_c$), \emph{style/culture} ($z_s$), and \emph{affective-functional relevance} ($z_a$).
We encourage $z_a$ to be less culture-discriminative through domain-adversarial training \citep{ganin2016dann}, then perform cross-cultural recommendation by aligning user preference distributions and target-culture candidate distributions with entropic OT \citep{cuturi2013sinkhorn}.
To refine ambiguous regions of the latent space, we attach a PAL loop \citep{settles2009active} that selects uncertain examples for expert judgment and converts them into pairwise constraints.

This paper makes three contributions.
First, we present a structured downstream recommendation framework that combines foundation embeddings with disentangled latent variables, domain-adversarial alignment, and OT ranking.
Second, we propose a unified evaluation protocol for cross-cultural recommendation that foregrounds serendipity \citep{zhang2012auralist}, cultural calibration, and disentanglement proxies alongside ranking quality.
Third, we outline a pilot human-in-the-loop setting in which small amounts of expert feedback can be injected back into the recommendation model.

\section{Related Work}
\subsection{Foundation Audio Models and Cross-Cultural MIR}
Music foundation models such as MERT have become strong generic front ends for MIR \citep{li2023mert}.
Yet their training data remain unevenly distributed across musical traditions, and performance on non-Western repertoires can lag without targeted adaptation.
CultureMERT addresses this issue through continual pre-training on a multi-cultural mixture and reports gains on non-Western auto-tagging tasks \citep{kanatas2025culturemert}.
In parallel, GlobalMood highlights the need for cross-cultural benchmarks and culturally grounded annotation protocols in music emotion research \citep{lee2025globalmood}.

Our work is complementary rather than competitive with these efforts.
We do not propose a new foundation model.
Instead, we ask what kind of \emph{downstream structure} is needed once strong audio embeddings are available but the target task is cross-cultural recommendation rather than tagging.

\subsection{Disentanglement and Invariance}
Disentangled latent-variable models offer one route for separating factors of variation in learned representations \citep{higgins2017betavae,kim2018factorvae}.
For recommendation, the relevant issue is not abstract factor disentanglement in isolation, but the separation of stylistic-cultural variation from the affective or functional information that may support meaningful cross-cultural matches.
We therefore use a pragmatic, task-driven notion of factorization rather than claiming perfect semantic disentanglement.

Domain-adversarial training encourages representations that are informative for the main task but less discriminative for nuisance domains \citep{ganin2016dann}.
In our setting, the nuisance domain corresponds to musical culture.
This provides a direct mechanism for shaping the latent subspace used for cross-cultural alignment.

\subsection{Serendipity, Recommendation, and Human Feedback}
Music recommenders should support discovery, not only accuracy.
Serendipity-oriented recommendation emphasizes the combination of relevance and unexpectedness \citep{zhang2012auralist}.
This is especially important in cross-cultural settings, where users may value recommendations that are unfamiliar in style but still emotionally or functionally resonant.

Active learning provides a natural framework for acquiring expensive human labels where they matter most \citep{settles2009active}.
Our PAL formulation adapts this logic to cross-cultural recommendation:
instead of treating human feedback as a large static annotation phase, we target boundary cases that are especially informative for latent-space repair.

\section{Method}
\subsection{Problem Setup}
Let $e_i \in \mathbb{R}^d$ denote a fixed audio embedding for track $i$, extracted from a foundation model such as CultureMERT.
Each track also has a culture label $c_i \in \mathcal{C}$.
User histories are represented as weighted sets of interacted tracks.
Given a user $u$ and a target culture $c^\star$, the goal is to recommend tracks from culture $c^\star$ that remain relevant to the user while increasing cross-cultural discovery.

We assume that raw foundation embeddings alone are insufficient because they entangle content, style, and shared affective-functional signals.
We therefore learn a downstream latent representation:
\begin{equation}
e_i \mapsto (z_{c,i}, z_{s,i}, z_{a,i}),
\end{equation}
where $z_c$ captures content structure, $z_s$ captures style/cultural variation, and $z_a$ captures relatively culture-agnostic affective-functional relevance.

\subsection{Disentangled Downstream Representation}
We use a variational encoder-decoder architecture to produce the three latent factors.
The overall objective combines reconstruction, KL regularization, consistency, and optional pairwise constraints:
\begin{equation}
\mathcal{L} =
\mathcal{L}_{\text{recon}}
+ \beta \mathcal{L}_{\text{KL}}
+ \lambda_{\text{sim}} \mathcal{L}_{\text{sim}}
+ \lambda_{\text{pair}} \mathcal{L}_{\text{pair}}
+ \lambda_{\text{dom}} \mathcal{L}_{\text{dom}}.
\end{equation}

Reconstruction preserves enough information to recover the input embedding.
KL regularization stabilizes the latent geometry.
The similarity term uses augmentation-based consistency inspired by contrastive learning \citep{oord2018cpc} so that content information remains stable under mild acoustic perturbations.
Pairwise constraints, when available, pull semantically similar items together and push dissimilar items apart.

\subsection{Domain-Adversarial Alignment}
To reduce cultural leakage in the latent variables used for cross-cultural matching, we attach a culture discriminator to $z_a$ and optimize it adversarially \citep{ganin2016dann}.
This is not intended to remove all cultural structure from the system.
Instead, it aims to reserve culture-specific variation primarily for $z_s$ while keeping $z_a$ more suitable for cross-cultural comparison.

\subsection{Optimal Transport Recommendation}
Given a user history $H_u$, we compute a source distribution in the $z_a$ space over the user's interacted items.
For a target culture $c^\star$, we compute a target distribution over candidate tracks from that culture.
We then solve an entropically regularized OT problem \citep{cuturi2013sinkhorn}:
\begin{equation}
\pi^\star = \arg \min_{\pi \in \Pi(\mu_u,\nu_{c^\star})}
\langle \pi, M \rangle - \varepsilon H(\pi),
\end{equation}
where $M$ is the pairwise distance matrix in $z_a$.
Candidate scores are derived from the resulting transport plan.
This shifts recommendation from pointwise nearest-neighbor matching to distributional alignment.

\subsection{Participatory Active Learning}
Model uncertainty is used to select examples near latent decision boundaries.
Experts inspect these items and provide pairwise judgments or concept-level annotations.
These judgments are converted into constraints and injected back into training.
In the present draft, we treat this as a pilot mechanism rather than a large-scale annotation program.

\section{Experimental Setup}
\subsection{Dataset Pipeline}
Our current data pipeline is built from public audio sources spanning four cultural domains:
West, India, Turkey, and China.
The internal dataset version used in this draft contains 1{,}600 tracks, balanced across the four domains, together with foundation embeddings, weak interactions, and fixed train/validation/test splits.
This paper draft uses the real dataset structure but \emph{synthetic placeholder outcomes} for the quantitative sections below.

\subsection{Baselines}
We compare against four baselines:
\begin{enumerate}
\item CultureMERT cosine retrieval.
\item CultureMERT + shallow MLP ranking head.
\item DCAS without OT.
\item DCAS without domain-adversarial alignment.
\end{enumerate}
The full model additionally includes pairwise-constraint support and the PAL feedback pathway.

\subsection{Metrics}
We report ranking quality (nDCG@10 and Recall@10), serendipity, and cultural calibration KL.
We also report disentanglement proxy metrics (MIG, DCI, SAP) for the latent variables.
Serendipity is operationalized as the product of unexpectedness in $z_s$ and relevance in $z_a$, following the intuition that cross-cultural discovery should be both unfamiliar and meaningful.

\section{Results}
\subsection{Main Recommendation Results}
Table~\ref{tab:main} shows the primary draft comparison.
The full model improves over embedding-only retrieval in both ranking and serendipity while also reducing calibration error.
Although these values are placeholders, they represent the intended empirical pattern to be validated in the final paper.

\begin{table*}[t]
\centering
\caption{Main recommendation results.
All values marked with $\dagger$ are synthetic placeholders for internal drafting only. Higher is better except Calibration KL.}
\label{tab:main}
\begin{tabular}{lcccc}
\toprule
Model & nDCG@10 & Recall@10 & Serendipity & Calibration KL \\
\midrule
CultureMERT cosine & 0.401$^\dagger$ & 0.512$^\dagger$ & 0.171$^\dagger$ & 0.91$^\dagger$ \\
CultureMERT + MLP & 0.417$^\dagger$ & 0.528$^\dagger$ & 0.180$^\dagger$ & 0.87$^\dagger$ \\
DCAS w/o OT & 0.425$^\dagger$ & 0.535$^\dagger$ & 0.218$^\dagger$ & 0.82$^\dagger$ \\
DCAS w/o domain adv. & 0.432$^\dagger$ & 0.543$^\dagger$ & 0.231$^\dagger$ & 0.80$^\dagger$ \\
\textbf{DCAS full} & \textbf{0.439}$^\dagger$ & \textbf{0.547}$^\dagger$ & \textbf{0.252}$^\dagger$ & \textbf{0.69}$^\dagger$ \\
\bottomrule
\end{tabular}
\end{table*}

This pattern supports the hypothesis that simple embedding retrieval is not enough for cross-cultural recommendation.
The OT component contributes to higher serendipity, while the domain-adversarial objective helps reduce the tendency to match items based on culture-specific style alone.

\subsection{Disentanglement and Ablation}
Table~\ref{tab:ablation} presents placeholder ablations.
The intended conclusion is not that latent factors are perfectly disentangled, but that the structured inductive bias improves the geometry of the downstream recommendation space enough to matter behaviorally.

\begin{table}[t]
\centering
\caption{Ablation study with synthetic placeholder values.}
\label{tab:ablation}
\begin{tabular}{lcc}
\toprule
Variant & Serendipity & Calibration KL \\
\midrule
Full model & 0.252$^\dagger$ & 0.69$^\dagger$ \\
No OT & 0.218$^\dagger$ & 0.82$^\dagger$ \\
No domain adv. & 0.231$^\dagger$ & 0.80$^\dagger$ \\
No constraints & 0.240$^\dagger$ & 0.74$^\dagger$ \\
\bottomrule
\end{tabular}
\end{table}

The latent evaluation protocol yields placeholder proxy scores of MIG $=0.024^\dagger$, DCI disentanglement $=0.229^\dagger$, and SAP $=17.3^\dagger$ for the best-performing run.
These numbers should be interpreted only as slots for the final experimental narrative.

\subsection{Participatory Feedback Pilot}
Finally, we simulate the structure of a small PAL experiment with one round of expert feedback.
The intended empirical claim is modest:
small but high-value human input should improve ambiguous cross-cultural boundaries more efficiently than uniform annotation.

\begin{table}[t]
\centering
\caption{Illustrative PAL pilot results with synthetic placeholder values.}
\label{tab:pal}
\begin{tabular}{lcc}
\toprule
Condition & Serendipity & Calibration KL \\
\midrule
Before PAL & 0.252$^\dagger$ & 0.69$^\dagger$ \\
After simulated PAL & 0.259$^\dagger$ & 0.66$^\dagger$ \\
After real PAL pilot & 0.268$^\dagger$ & 0.62$^\dagger$ \\
\bottomrule
\end{tabular}
\end{table}

\section{Discussion}
This work should not be read as a claim that disentanglement alone solves cross-cultural recommendation.
Instead, the key argument is narrower.
Cross-cultural recommendation fails in part because systems often conflate unfamiliar style with irrelevance.
A structured downstream representation provides a useful inductive bias for separating these signals.

Likewise, the purpose of OT is not mathematical novelty for its own sake.
Its practical value is that it aligns preference distributions rather than individual nearest neighbors, reducing collapse toward culturally dominant local matches.
Finally, PAL matters because cross-cultural boundary cases are hard to learn from weak logs alone.
Some distinctions require direct human judgment.

\section{Ethics Statement}
This project studies recommendation across musical cultures and therefore raises questions of representation, annotation power, and cultural reduction.
We do not claim that a three-factor latent model fully captures musical meaning across cultures.
Rather, we use it as a limited engineering approximation for recommendation.
Any deployment or publication of such a system should disclose dataset composition, annotation provenance, and limitations in cultural coverage.
The PAL component is motivated by the need to avoid treating culturally grounded concepts as fixed labels imposed from outside the communities in question.

\section{Conclusion}
We presented a proof-of-concept framework for cross-cultural music recommendation that combines foundation embeddings with disentangled downstream representation learning, domain-adversarial shaping, OT alignment, and a participatory feedback pathway.
The central claim of this draft is not that large audio models are unnecessary, but that they benefit from additional downstream structure when the target problem is culturally sensitive recommendation rather than generic retrieval.
While the present manuscript uses synthetic placeholder outcomes for internal writing purposes, the paper structure, hypotheses, and evaluation design are intended to support a final ISMIR submission once the remaining experiments are complete.

\bibliographystyle{plainnat}
\bibliography{refs}

\end{document}
